%! TEX root = ../am1-19-20.tex

\chapter{Examen 2018--2019}

\indent\indent 1. Calculați integrala:
\[
  I = \int_{-\infty}^\infty \frac{e^{2x}}{(1 + e^{2x})^2} dx,
\]
folosind integralele gamma și/sau beta.

\vspace{0.5cm}

2. Calculați lucrul mecanic efectuat de forța $ \vec{F} = xy^2 \vec{i} + 2x^2y \vec{j} $
asupra unei particule care circulă în linie dreaptă între punctele $ O(0, 0), A(2, 2), B(2, 4) $
și înapoi la $ O $ (în sens direct). Reprezentare grafică.

\emph{Indicație:} Lucrul mecanic se calculează ca integrală curbilinie
de speța a doua.

\vspace{0.5cm}

3. Să se calculeze $ \ds\iint_D 1 + \sqrt{x^2 + y^2} dxdy $, unde
$ D = \{ (x, y) \in \RR^2 \mid x^2 + y^2 - y \leq 0, x \geq 0 \} $.

\emph{Indicație:} Se folosesc coordonate polare (translatate! cercul nu este
centrat în origine) sau se descrie domeniul ca unul intergrafic.

\vspace{0.5cm}

4. Calculați volumul corpului cuprins între suprafețele de ecuație $ x^2 + y^2 = -z $
și $ z = y $. Reprezentare grafică.

\emph{Indicație:} Se calculează integrală triplă de la o suprafață la cealaltă.

\vspace{0.5cm}

5. Se consideră suprafața $ S $ de ecuație $ 1 - z = (x - 2)^2 + y^2 $ cu
bordul orientat $ \partial S : 1 - z = (x - 2)^2 + y^2$, cu $ z = 0 $.

Folosind formula lui Stokes, calculați integrala:
\[
  \int_{\partial S} (x^2 + y^3) dx + dy + zdz,
\]
în care orientarea este în sensul de creștere a lui $ z $.

%%% Local Variables:
%%% mode: latex
%%% TeX-master: "../am1-19-20"
%%% End:
